\section*{ВВЕДЕНИЕ}
\addcontentsline{toc}{section}{ВВЕДЕНИЕ}

В современных организациях значительная часть деятельности связана со сбором, регистрацией, хранением и обработкой данных пользователей. Информация о клиентах, сотрудниках, обращениях и выполняемых операциях является важнейшим ресурсом, от качества управления которым зависит эффективность работы предприятия. Использование бумажного документооборота и разрозненных электронных таблиц приводит к дублированию информации, потере актуальности данных, затрудняет поиск необходимых сведений и не позволяет обеспечить прозрачность процессов и контроль действий пользователей.

Особую актуальность задача автоматизации регистрации данных пользователей приобретает в организациях сервисного обслуживания, где одновременно ведётся работа с большим количеством клиентов и связанных с ними объектов учёта. На примере автосервисных организаций возникает необходимость не только хранить сведения о клиентах и автомобилях, но и обеспечивать регистрацию заявок на обслуживание, контролировать их исполнение, распределять обязанности между сотрудниками и фиксировать все изменения, происходящие в системе.

На сегодняшний день существует множество программных решений, предназначенных для автоматизации взаимоотношений с клиентами и учёта бизнес-процессов. Крупные CRM-системы обладают широким функционалом, однако зачастую отличаются высокой стоимостью внедрения и сопровождения, сложностью настройки и содержат избыточные возможности, невостребованные малыми и средними организациями. Использование универсальных систем учёта не всегда позволяет учитывать особенности конкретной предметной области и требует дополнительной адаптации. Применение электронных таблиц, несмотря на простоту использования и отсутствие затрат на приобретение программного обеспечения, не обеспечивает разграничение прав доступа, ведение журнала изменений и возможность одновременной работы нескольких пользователей.

Таким образом, разработка программно-информационной системы регистрации данных пользователей, ориентированной на решение задач конкретной организации, является актуальной задачей. Использование современных веб-технологий позволяет обеспечить доступ к системе из любого места при наличии подключения к сети Интернет, а также упростить сопровождение и дальнейшее развитие программного продукта.

Целью настоящей работы является разработка программно-информационной системы регистрации данных пользователей в виде веб-приложения, обеспечивающей учёт пользователей, клиентов, автомобилей и заявок на обслуживание, разграничение прав доступа и ведение журнала изменений действий пользователей.

Для достижения поставленной цели необходимо решить следующие задачи:
\begin{itemize}
	\item провести анализ предметной области и обосновать необходимость автоматизации процессов регистрации и учёта данных пользователей;
	\item спроектировать архитектуру программной системы и структуру базы данных;
	\item разработать подсистемы управления пользователями, клиентами, автомобилями и заявками с возможностью выполнения операций создания, просмотра, редактирования и удаления данных;
	\item реализовать механизмы аутентификации и авторизации пользователей на основе ролевой модели;
	\item реализовать регламентированный жизненный цикл заявок и механизм назначения исполнителей;
	\item обеспечить разграничение прав доступа на уровне приложения и базы данных;
	\item реализовать автоматическое ведение журнала изменений ключевых сущностей системы;
	\item провести функциональное тестирование разработанной программной системы.
\end{itemize}

Предметом исследования являются методы и средства регистрации и обработки данных пользователей, механизмы аутентификации и авторизации, модели разграничения прав доступа, алгоритмы управления жизненным циклом заявок, а также средства ведения аудита действий пользователей в веб-приложениях.

Объектом разработки является программно-информационная система регистрации данных пользователей, реализованная в виде CRM-системы для автоматизации деятельности автосервисной организации.

Структура и объём работы. Выпускная квалификационная работа состоит из введения, четырёх разделов основной части, заключения, списка использованных источников и двух приложений. Текст выпускной квалификационной работы составляет \formbytotal{lastpage}{страниц}{у}{ы}{}.

Во введении сформулирована цель работы, поставлены задачи разработки, определены объект и предмет исследования, описана структура работы и приведено краткое содержание её разделов.

В первом разделе на стадии анализа предметной области рассматриваются особенности деятельности организации сервисного обслуживания, выявляются процессы, подлежащие автоматизации, анализируются существующие программные решения и обосновывается необходимость разработки собственной программно-информационной системы.

Во втором разделе на стадии технического задания формулируются основание для разработки, назначение и цели системы, требования к функциональности, безопасности, программному и аппаратному обеспечению, а также моделируются варианты использования для различных категорий пользователей.

В третьем разделе на стадии технического проектирования представлены архитектура программной системы, структура базы данных, описание основных сущностей и взаимосвязей между ними, а также обоснование выбора используемых технологий и проектных решений.

В четвёртом разделе приводится описание реализации программной системы, рассматриваются ключевые механизмы аутентификации, авторизации, управления заявками и ведения журнала изменений, представлены результаты функционального тестирования и рекомендации по развёртыванию системы.

В заключении излагаются основные результаты работы, полученные в ходе разработки.

В приложении А представлен графический материал.

В приложении Б представлены фрагменты исходного кода программы.
